\subsection{Uniformity over cuts of open chains}

\begin{corollary}[Arbitrary end cuts of the finite chains]
\label{cor:chain-arbitrary-cuts}
For the half-filled open XX chain of even length $N$, take the site
bipartition $\{1,\ldots,\ell\}\mid\{\ell+1,\ldots,N\}$ and put
$m=\min(\ell,N-\ell)$. As $m\to\infty$, uniformly in $N$ and $\ell$,
\[
 S=\tfrac16\log m+O(1),\qquad V=\tfrac16\log m+O(1),\qquad
 \mathcal E_d=\frac{3(\log d)^2}{2\log m}+O_d((\log m)^{-2}).
\]
For the open critical Ising chain of any length $L$ cut after $k$ spins,
put $m=\min(k,L-k)$. The same statement holds with
$S,V=\tfrac1{12}\log m+O(1)$ and
$\mathcal E_d=3(\log d)^2/\log m+O_d((\log m)^{-2})$.
\end{corollary}

\begin{proof}
By reflection it suffices to treat the XX prefix cut with $\ell=m$.
Write $M=N+1$. For $1\le x\le m<y\le N$, the exact cross-covariance
entry, up to the sign of the filled standing-wave half, is
\[
 \frac{\sin(\pi(x-y)/2)}{2M\sin(\pi(x-y)/(2M))}
 -\frac{\sin(\pi(x+y)/2)}{2M\sin(\pi(x+y)/(2M))}.
\]
This follows by expanding the product of standing waves and summing
the first $N/2$ cosines. In both terms replace the reciprocal factor
$(2M\sin(\pi(x\pm y)/(2M)))^{-1}$ by $(\pi(x\pm y))^{-1}$.
The error in trace norm is $O(1)$ uniformly in $m,N$. Indeed the
difference is $M^{-1}f((x\pm y)/M)$ with
$f(z)=(2\sin(\pi z/2))^{-1}-(\pi z)^{-1}$ analytic for $|z|<2$.
For a matrix $M^{-1}(a_x+b_y)^q$, the binomial rank-one expansion
bounds its trace norm by
$M^{-1}\sqrt{m(N-m)}(\max_x|a_x|+\max_y|b_y|)^q$.
Here the sum of maxima is at most $3/2$ and the prefactor is at most
$1/2$. The Taylor series of $f$ is absolutely summable on that
smaller radius. Multiplication by either sine gives a linear
combination of two diagonal-unitary left and right multipliers,
so the same bound applies to the covariance remainder.

We use a uniform bound for the rectangular matrix with entries
$(j+k+a)^{-1}$ for $0\le j<p$ and $p\le k<q$, with fixed $a>0$ and
$1\le p\le q$. Its Laplace rank-one representation bounds the trace
norm by
\[
 \int_0^\infty e^{-(p+a)t}
 \frac{\sqrt{\min\{p,(1-e^{-2t})^{-1}\}}}{\sqrt{1-e^{-2t}}}\,dt.
\]
On $0<t\le1/p$ the integrand is bounded by $C\sqrt{p/t}$.
On $t\ge1/p$ it is bounded by $C e^{-pt}(1+t^{-1})$. Both
integrals are bounded independently of $p,q$. Taking $a=2$ shows that
the reflected principal matrix with entries
$\sin(\pi(x+y)/2)/(\pi(x+y))$ has trace norm $O(1)$, using the same
sine phase decomposition.

Reverse the order on the left by setting $u=m+1-x$ and set $v=y-m$
on the right. The remaining principal matrix has entries
$\sin(\pi(u+v-1)/2)/(\pi(u+v-1))$. Its parity blocks, after removing
alternating signs, are rectangular Hilbert matrices scaled by
$(2\pi)^{-1}$, with $a=1/2,3/2$ and row counts
$m_1=\lceil m/2\rceil$, $m_2=\lfloor m/2\rfloor$. Each column count
is at least its row count. The rectangular bound just proved gives
trace-norm error $O(1)$ when the columns beyond the leading square
are deleted with zero padding. Lemma~\ref{lem:hilbert-mesh} and the
singular-value perturbation inequality therefore give the comparison
\eqref{eq:chain-main-mesh} with $L_i=\log(2m_i+1)=\log m+O(1)$.
All constants in the comparison are independent of the cut. The
proof of Theorem~\ref{thm:xx-full-law} applies with $W=\log m$.

For Ising, reflect so that $k\le L/2$ and take the XX chain on $2L$
sites with a cut after $2k$ sites. In the Majorana decomposition of
Corollary~\ref{cor:ising-fourier}, each path then restricts to its
first $2k$ Majoranas, the first $k$ Ising spins. The XX mode list is
again the multiset union of two Ising lists. The XX result at
$W=\log(2k)$ and exact doubling prove the claimed Ising estimates
uniformly in $L,k$.
\end{proof}
